\section{未来方向与开放问题}

\subsection{数据是核心瓶颈}

\begin{figure}[H]
\centering
\includegraphics[width=0.95\textwidth]{slides-images/slide-083.jpg}
\caption{当前研究方向总结\protect\footnotemark}
\end{figure}
\footnotetext{Slides 第84页。}

Lambert 认为对齐研究最大的瓶颈不是算法，而是\textbf{数据}：

\begin{itemize}
    \item 学术界几乎\textbf{试遍了所有公开数据集}——但选择仍然非常有限
    \item UltraFeedback 已经使用了半年多，按照行业标准已经"过时"
    \item 新数据集的创建需要大量人力和资源，但社区缺乏可持续的数据收集机制
\end{itemize}

\subsection{五大研究方向}

Lambert 列出了他最关注的五个方向：

\begin{enumerate}
    \item \textbf{数据创新}：需要新的偏好数据集，特别是能够提升特定能力（如事实性、推理）的数据

    \item \textbf{DPO 方法的变体}：包括去除参考模型、修改损失函数、使用单侧偏好（如 KTO 方法）等

    \item \textbf{小模型对齐}：在 7B 甚至更小的模型上进行对齐研究——这是学术界可以发力的方向，因为大公司都在竞争更大的模型

    \item \textbf{更好的评估工具}：需要更具体、更有针对性的评估基准，而非通用排行榜

    \item \textbf{个性化对齐}：训练适合个人需求的模型，而非一个通用模型服务所有人
\end{enumerate}

\begin{importantbox}{超越二元偏好的可能性}
Q\&A 环节中讨论了是否存在比成对偏好更好的反馈形式。Lambert 列举了几个方向：
\begin{itemize}
    \item \textbf{KTO}（Stanford）：使用单侧偏好（yes/no）而非成对比较
    \item \textbf{K-wise 偏好}（Starling 模型）：同时比较 5--9 个回复并排序
    \item \textbf{细粒度偏好}：为每个回复标注多个维度（简洁性、有用性、诚实性等）
    \item \textbf{社会选择理论}：整个社会选择（social choice）领域的知识还未被充分引入
\end{itemize}
\end{importantbox}

\subsection{Reward Hacking 是永恒的挑战}

在 Q\&A 环节中，关于 reward hacking 的讨论引发了深入思考：

\begin{warningbox}{Reward Hacking 不可完全避免}
当你有一个强大的优化器（如 PPO）和一个不完美的奖励表示时，优化器\textbf{总会}找到奖励表示的漏洞。Lambert 举了一个有趣的例子：如果训练不加约束，模型可能会对所有问题都回答"JavaScript"——因为这恰好能在某些评估指标上获得高分。KL 散度约束就是用来缓解这个问题的"刹车"，但不可能完全消除 reward hacking。
\end{warningbox}

\subsection{超越人类水平的对齐}

另一个重要问题是：如果对齐依赖于人类反馈，如何让模型超越人类水平？Lambert 认为\textbf{搜索（search）}将是关键：

\begin{itemize}
    \item 搜索本质上是 RL 中的\textbf{探索（exploration）}
    \item 语言模型可以通过搜索生成新的合成数据
    \item 人类的作用将从"提供所有答案"转变为"验证和纠正搜索结果"
    \item 这可能是 OpenAI 的 Q* 项目试图解决的问题
\end{itemize}

\subsection{本章小结}

对齐研究正处于快速发展但充满不确定性的阶段。数据是最大的瓶颈，算法创新空间仍然广阔，评估工具亟需完善。未来的方向可能包括：在线方法的成熟、超越二元偏好的反馈机制、小模型对齐、个性化对齐，以及利用搜索突破人类水平的限制。

%% ========== 总结与延伸 ==========

